% Transcription — fonds Alexandre Grothendieck, shelfmark 123, batch 1 (pages 1-20).
% Established from the facsimile archives/batches/123/batch-01.pdf (a mirror of
% the folder's PDF; no cover sheet, so PDF page k = folder page k), per the
% skill transcribe-grothendieck.
%
% Pass: Opus 5.5 (claude-opus-5-5), 2026-09-24 — first pass, unchecked against the pages by a human.
%
% \dating{} carries the folder's own inventory date, "1969", verbatim; per the
% inventory that date belongs to the letter, the typescripts and notes being
% s.d. The group "119-121" (« Autour d'Esquisse d'un programme », 1968-1991)
% is not used for dating.
%
% Sixteen of the twenty pages are transcribed. Absent:
%   - pages 3, 7 and 14, covers bearing only a title in his hand; the titles
%     head the three \section*s that follow them (hand.md §5):
%     3 « Lissification des morphismes de Kostant », 7 « Lissification des
%     morphismes de Steinberg », 14 (buff card) « Résumé th. de conjugaison
%     Kostant + lettre à Kostant »;
%   - page 20, a blank buff sheet.
% The letter to Kostant announced on the cover of page 14 is not in this
% batch.
%
% The typescripts (pages 1-2 and 15-19). The brief's rule is that another
% person's typescript is summarised page by page and only his marks given in
% full. These two are judged HIS OWN typed texts, not someone else's, and
% are therefore transcribed in full like his manuscripts: both are in French,
% cite Kostant in the third person ("Grâce à Kostant", "Deligne nous
% fournit", "j'ignore dans quelle mesure"), are filed under his own covers,
% and are corrected throughout in his hand, which also inks in the gothic
% letters the machine lacked. If a later check shows otherwise they would
% have to be cut back to summaries. Typing slips overstruck with x at the
% machine are not recorded.
%
% What the batch is.
%   1-2    Typescript « Compléments sur transformation de Coxeter d'un système
%          de racines (Kostant) »: eigenvalues of w in W are roots of unity of
%          order dividing some m_i+1; order h iff w is Coxeter, eigenvectors
%          then regular; Cor. 1 the centraliser of a Coxeter element c is <c>
%          and N/<c> = (Z/hZ)^* (the typescript writes wcw^-1 = c^h where c^n
%          is meant — flagged, not corrected); x regular iff eigenvalue a
%          primitive h-th root; cyclic lines, card(W)/h of them.
%   4-6    Grothendieck-Springer for the Lie algebra over a base S: g, t, W,
%          J = t/W, Bor, Kill, the universal Borel and its Lie algebra L; the
%          cartesian square over g^rss, the morphism f: g -> J (Kostant),
%          g' = g x_J J', p: L -> g' projective and birational; Springer PM
%          n° 30; a list of Questions 0)-g) (flatness of f, regularity of J,
%          fibres, smooth points, orbital stratification). Page 6's questions
%          c)-e) are fast and read only in fragments.
%   8-13   The group version: G reductive over S, D = Bor(G), B universal
%          Borel, T, K = Kill(G), W, I = T/W, Steinberg's f: G -> I, G' =
%          G x_I T, p: B -> G'. Theorem (i)-(iv): p projective, f~ smooth,
%          B^reg -> G'^reg iso (Main Theorem + Steinberg), codimension
%          statements, Pic(B) = Pic(B^reg) = Pic(G'^reg) <- Pic(D), ample cone
%          = P^+; proof of (iv) via EGA IV 21.4.13 and a Lemma Pic(T) =
%          H^1(Γ,M); Remarks a)-d): cokernel of Pic(D) -> Pic(B) finite,
%          Pic(S) = Pic(G'), and the conclusion that no resolution of f
%          analogous to that of f' exists without touching G^reg.
%   15-19  Typescript « Résumé de quelques résultats de Kostant »: the
%          Jacobson-Morozov / Kostant theorem 1 (i)-(iv) (conditions a)-d) on
%          (x,e), conjugacy under Cent(u), torsors under Z^+, g(2)^*, the
%          parabolic P, U, L of a Dynkin element); then « Interprétations en
%          termes de filtrations du foncteur oubli »: classes c, schemes
%          ST_c, Ø_c, Par_c', Dyn_c, Un_c, Deligne's Un_c -> Ø_c, rational
%          points and quasi-split groups, Dynkin-Kostant parabolics, the
%          principal case. His margins: A_2, A_4 (injectivity of c -> c',
%          subregular unipotents giving Borels).
%
% Notation fixed for the batch: \mathfrak{g} for his gothic g (underlined when
% it is the Lie algebra scheme, pp. 4-6), \underline{t}, \mathfrak{J},
% \underline{\mathfrak{L}}, \mathfrak{B} for the round B of the universal
% Borel, \mathfrak{z} for the curly 3 (pp. 1, 18), \underline{\varnothing}
% for the typed Ø of the filtration schemes.
\documentclass[11pt,a4paper]{article}
% Le préambule commun à toutes les transcriptions du fonds Grothendieck.
%
% Il tient en une idée : **l'incertitude se compose, elle ne se résout pas.**
% Une transcription qui lisse un mot douteux a détruit la seule chose pour
% laquelle on la faisait — un lecteur doit pouvoir savoir, à chaque endroit,
% ce qui était sur le feuillet et ce qui a été ajouté. Les six macros qui
% suivent sont donc l'appareil critique, et non de la décoration.
%
% Les mêmes commandes sont reconnues par `scripts/render.mjs`, qui produit la
% vue de lecture du site. Une macro ajoutée ici doit l'être aussi là-bas,
% faute de quoi le rendu échouera bruyamment — ce qui est voulu : un rendu
% silencieusement faux est pire qu'un rendu absent.

\ProvidesPackage{grothendieck}[2026/08/08 Transcriptions du fonds Grothendieck]

\usepackage{amsmath,amssymb,amsthm}
\usepackage{mathrsfs}
\usepackage{stmaryrd}
% Les diagrammes commutatifs sont partout dans ce fonds : c'est la langue
% même de la géométrie algébrique de Grothendieck, et une transcription qui
% les rendrait en prose ne transcrirait rien.
\usepackage{tikz-cd}
% Français par défaut — c'est la langue des feuillets — mais l'anglais est
% chargé aussi : la traduction et le résumé sont des documents anglais, et la
% typographie française (espaces avant les deux-points) y serait une faute.
% Ces documents basculent par \selectlanguage{english} en tête de corps.
\usepackage[english,french]{babel}
\usepackage{fontspec}
\usepackage{geometry}
\geometry{a4paper,margin=25mm,marginparwidth=28mm}
\usepackage{marginnote}
\usepackage{xcolor}
% ulem, not soul. `soul`'s \st and \ul are built on a character-by-character
% scan that XeLaTeX's font machinery breaks: the document fails with an
% undefined control sequence rather than degrading. ulem does the same two jobs
% and is Unicode-engine safe. `normalem` stops it hijacking \emph, which the
% transcriptions use for Grothendieck's own emphasis.
\usepackage[normalem]{ulem}
\usepackage{hyperref}
\hypersetup{colorlinks=true,linkcolor=black,urlcolor=black,pdfborder={0 0 0}}

% --- Métadonnées du lot -----------------------------------------------------
% Reprises telles quelles par le script de rendu, qui les affiche en tête de
% la vue de lecture. Elles fixent ce que le fichier prétend couvrir : sans
% elles, une transcription flotte sans point d'attache dans un fonds de seize
% mille feuillets.

\newcommand{\folder}[1]{\gdef\grothFolder{#1}}
\newcommand{\batch}[1]{\gdef\grothBatch{#1}}
\newcommand{\leaves}[2]{\gdef\grothFirst{#1}\gdef\grothLast{#2}}
\newcommand{\foldertitle}[1]{\gdef\grothFoldertitle{#1}}
\gdef\grothFolder{?}\gdef\grothBatch{?}\gdef\grothFirst{?}\gdef\grothLast{?}\gdef\grothFoldertitle{}

% --- Le résumé --------------------------------------------------------------
%
% Il ouvre la lecture modernisée, sous le titre « Résumé », et il est composé
% en petit. Ce n'est pas un ornement : c'est le seul passage du document qui
% ne soit pas des mathématiques, et un lecteur doit pouvoir le reconnaître
% avant de l'avoir lu — puis le sauter s'il connaît déjà le sujet, ou s'y
% arrêter s'il ne le connaît pas. Le filet qui le suit marque la reprise du
% corps.

\newenvironment{resume}
  {\small\setlength{\parskip}{0.45em}}
  {\par\medskip\hrule\medskip}

% --- Mots-clés ---------------------------------------------------------------
%
% En anglais, en clôture du résumé de la lecture modernisée : le vocabulaire
% moderne sous lequel chercher ce que les feuillets construisent. Le manifeste
% du site (`npm run manifest`) les extrait pour étiqueter la cote entière —
% cette ligne est la source unique des tags, il n'y a pas de fichier à côté.
\newcommand{\keywords}[1]{\par\smallskip\noindent
  {\footnotesize\textsc{Keywords} --- \textit{#1}}\par}

% --- L'appareil critique ----------------------------------------------------

% Le numéro de feuillet, en marge. C'est l'ancre qui permet de régler un
% désaccord sur une formule en regardant *un* feuillet plutôt qu'une chemise
% de six cents. Le site s'en sert aussi pour tourner les pages du fac-similé
% quand on fait défiler la transcription.
\newcommand{\leaf}[1]{%
  \marginnote{\footnotesize\color{gray!70}\textsf{#1}}%
  \label{leaf:#1}\ignorespaces}

% Illisible : on ne devine pas. Un mot inventé qui se lit comme les autres est
% le pire résultat possible.
\newcommand{\ill}{\textcolor{red!60!black}{[\,\ldots\,]}}

% Lecture douteuse : le mot est donné, et signalé comme incertain.
\newcommand{\uncertain}[1]{\uline{#1}}

% Ajout de l'éditeur — une lettre manquante, un mot sous-entendu.
\newcommand{\add}[1]{\textcolor{blue!50!black}{[#1]}}

% Biffé par Grothendieck lui-même. Ce qu'il a rayé fait partie du document :
% une rature dit souvent où la pensée a changé de direction.
\newcommand{\struck}[1]{\sout{#1}}

% Note du transcripteur, sur l'état matériel du feuillet.
\newcommand{\note}[1]{\par\smallskip{\small\color{gray!60!black}\textit{[#1]}}\par\smallskip}

% Note marginale de Grothendieck, distincte du corps du texte.
\newcommand{\marginal}[1]{\par\smallskip\noindent\hspace{1em}%
  {\small\color{gray!60!black}#1}\par\smallskip}

% --- Titre ------------------------------------------------------------------

\AtBeginDocument{%
  \begin{center}
    {\large\bfseries Fonds Alexandre Grothendieck --- cote n\textsuperscript{o}~\grothFolder}\\[2pt]
    {\itshape\grothFoldertitle}\\[4pt]
    {\small feuillets \grothFirst--\grothLast\ (lot \grothBatch)}\\[2pt]
    {\footnotesize\color{gray!60!black}Transcription établie d'après le fac-similé numérisé
     par l'Université de Montpellier.\\
     Les lectures incertaines sont soulignées, les illisibles notées [\,\ldots\,],
     les ajouts entre crochets.}
  \end{center}
  \vspace{1em}}

\endinput


\folder{123}
\batch{1}
\pages{1}{20}
\dating{1969}
\watermark{Édition de démonstration}
\foldertitle{[Autour de Kostant] : tapuscrits (s.d.), notes manuscrites (s.d.), lettre (1969).}

\begin{document}

\section*{Compléments sur la transformation de Coxeter (d'après Kostant)}

\note{pages 1 et 2 : deux feuillets dactylographiés, corrigés et annotés à
l'encre de sa main. Le titre ci-dessus est celui de l'éditeur ; le titre
dactylographié ouvre la page 1. Les fautes de frappe corrigées à la machine
(lettres ou mots recouverts de \emph{x}) ne sont pas relevées, sauf quand le
mot recouvert se lit et que la correction change le sens.}

\page{1}
Compléments sur transformation de Coxeter d'un système de racines
\marginal{irréductible}
(Kostant).
\marginal{The principal 3-dimensional subgroup …\ \ Th.\ 9.2}
\note{les deux ajouts sont à l'encre : \emph{irréductible}, entouré, est
appelé par un trait devant « (Kostant) » ; la ligne anglaise est écrite sous
le titre et renvoie, semble-t-il, au titre de l'article de Kostant et à son
théorème 9.2.}

a)\ Soit $w\in W$. Alors toute valeur propre de $w$ est une racine de $1$
d'ordre divisant un des entiers $m_i+1$ (où les $m_i$ sont les exposants
caractéristiques), donc d'ordre $\leqslant h$. Pour qu'il y en ait une qui
soit d'ordre $h$, il faut et suffit que $w$ soit une transformation de
Coxeter, et \add{alors} les éléments non nuls du sous-espace propre
correspondant
\marginal{\uncertain{(qui est de dim.\ 1)}}
sont des éléments réguliers.
\note{« alors » est ajouté à l'encre au-dessus de la ligne ; la note marginale
gauche est appelée par un signe d'insertion après « correspondant ».}

\textbf{Corollaire 1}\,: Soit $\mu$ le sous-groupe de $W$ engendré par un
élément de Coxeter $c$, donc $\mu\simeq \mathbf{Z}/h\mathbf{Z}$. Alors $\mu$
est égal au centralisateur de $c$, donc de $\mu$, et si $\nu$ est le
normalisateur de $\mu$, l'homomorphisme
\[
  \nu/\mu \longrightarrow \operatorname{Aut}(\mu)\simeq(\mathbf{Z}/h\mathbf{Z})^{*}
\]
est un isomorphisme\,; en d'autres termes, pour tout entier $n$ premier à $h$,
il existe un $w\in W$ tel que $wcw^{-1}=c^{h}$, et $w$ est déterminé
$\operatorname{mod}\mu$.
\note{le tapuscrit porte bien $c^{h}$ ; le sens, et la phrase qui suit
(« si $n$ est un entier premier à $h$, alors $c^{n}$ est conjugué de $c$ »),
demandent $c^{n}$. Non corrigé ici. Les lettres $\mu$, $\nu$ et les symboles
$\simeq$, $\in$ sont ajoutés à la main dans les blancs de la frappe.}
\marginal{Quelle extension de $(\mathbf{Z}/h\mathbf{Z})^{*}$ par
$\mathbf{Z}/h\mathbf{Z}$ obtient-on\,?}
\marginal{\uncertain{échanger avec cor.\ 3\,?}}
\note{la première note est dans la marge droite, en face de l'isomorphisme ;
la seconde dans la marge gauche, le long d'un trait vertical qui tient tout
l'énoncé du corollaire 1.}

Soit en effet $D$ \add{le sous-espace} propre sous $c$ correspondant à une
racine primitive $h$.ème de $1$. Alors $\dim D=1$ (prop.~30), et
$\mathfrak{z}=\operatorname{Centr}(\mu)$ normalise $D$, d'où un homomorphisme
$\mathfrak{z}\to\operatorname{Aut}(D)=\mathbf{C}^{*}$, dont l'image est un
sous-groupe fini de $\mathbf{C}^{*}$, contenant l'image (isomorphe) de $\mu$,
donc cyclique d'ordre un multiple de $h$. Grâce à Kostant, son ordre est égal
à $h$, et on a donc trouvé une rétraction $\mathfrak{z}\to\mu$, il faut
prouver que son noyau est réduit à $1$. Or si $w$ est dans ce noyau, il
\struck{induit} centralise $D$, et comme $D$ est régulière par Kostant, cela
implique $w=1$. On a donc prouvé $\mathfrak{z}=\mu$. Reste à prouver que si
$n$ est un entier premier à $h$, alors $c^{n}$ est conjugué de $c$, ce qui
résulte également de Kostant.
\note{« le sous-espace » est écrit à la main au-dessus d'un mot recouvert à
la machine. La lettre gothique notée ici $\mathfrak{z}$ est une figure de
3 bouclée, ajoutée à la main.}

\textbf{Corollaire}. Soit $c$ un élément de Coxeter, et $\lambda$ une valeur
propre de $c$, correspondant à un vecteur propre $x$. Pour que $x$ soit
régulier, il faut et il suffit que $\lambda$ soit une racine primitive
$h$.ème de $1$.
\note{un trait vertical, dans la marge gauche, tient cet énoncé.}

La suffisance a déjà été énoncée avec le th.\ de Kostant. Pour la nécessité,
soit $h'\leqslant h$ l'ordre de la valeur propre, alors $c^{h'}$ centralise
$x$, donc, $x$ étant régulier, on a $c^{h'}=1$, donc $h'=h$, cqfd.

\textbf{Corollaire}. \struck{Soit D une} \struck{Considérons} droite
régulière\add{s} dans $V\otimes\mathbf{C}$ telle qu'il existe un élé-
\note{deux débuts se superposent à la frappe, « Soit D une » et, au-dessus,
« Considérons », tous deux recouverts de caractères ; la phrase reprend à
« droite régulières ». La lecture voulue est vraisemblablement « Soit $D$ une
droite régulière ».}

\page{2}
ment de Coxeter \add{$c$} de $W$ qui la normalise. Alors le normalisateur de
$D$ dans $W$ est le groupe engendré par $c$, et il existe une unique élément
de Coxeter de $W$ normalisant $D$ et y ayant pour valeur propre une racine
$h$.ème primitive donnée de $1$. Les droites régulières $D$ ayant la propriété
précédente (appelées « droites cycliques » par Kostant) sont conjuguées sous
$W$, leur nombre est égal à $\operatorname{card}(W)/h$.
\note{le $c$ est ajouté à l'encre ; « $h$.ème » est frappé au-dessus de la
ligne.}

NB Dans l'ensemble des droites cycliques, le groupe
$(\mathbf{Z}/h\mathbf{Z})^{*}$ opère en commutant à l'action de $W$
librement\,; deux droites cycliques sont conjuguées sous l'action de ce groupe
si et seulement si elles ont même stabilisateur dans $W$.
\note{le reste du feuillet est blanc.}

\section*{Lissification des morphismes de Kostant}

\note{titre de sa main, en haut à droite du feuillet 3, qui ne porte rien
d'autre : c'est une chemise, et elle ne reçoit pas de \texttt{page}. Le
premier mot est d'une lecture sûre au feuillet 7, où il revient.}

\page{4}
\note{notations : $\underline{\mathfrak{g}}$ transcrit le $\mathfrak{g}$
gothique souligné de l'algèbre de Lie de $G$, $\mathfrak{J}$ le $J$ gothique
de l'espace des invariants, $\underline{t}$ l'algèbre de Lie du tore,
$\underline{\mathfrak{L}}$ un $L$ ronde soulignée, $\mathfrak{B}$ un $B$
rond. La feuille est une liste de notations, symbole à gauche, définition à
droite.}

\begin{itemize}
\item[$G$] schéma en groupes semi-simple \emph{adjoint} sur $S$.
\item[$\underline{\mathfrak{g}}$] algèbre de Lie \uncertain{comme Module} sur
$\mathrm{Sch}_{/S}$, donc schéma vectoriel sur $S$
\item[$\underline{\mathfrak{g}}^{\mathrm{rss}}$] sous-schéma ouvert des
éléments « réguliers \add{semi-}simples » i.e.\ de nullité minimum
(« réguliers » dans SGA 3 XIII)
\item[$T$] « le » tore maximal d'un couple de Killing (tore sur $S$)
\item[$\underline{t}$] son algèbre de Lie, regardée comme schéma vectoriel
sur $S$
\item[$W$] « le » groupe de Weyl, opérant sur $T$, donc sur $\underline{t}$.
\item[$\mathfrak{J}=\underline{t}/W$\,;] $\underline{t}$ est parfois noté
$\mathfrak{J}'$, donc $\mathfrak{J}'\to\mathfrak{J}$ revêtement (ramifié).
\item[$\underline{t}^{\mathrm{rss}}$] ouvert de $\underline{t}$ des éléments
« réguliers \add{semi-}simples ». Il est stable sous $W$.
\item[$\mathfrak{J}^{\mathrm{rss}}=\underline{t}^{\mathrm{rss}}/W$\,;]
$\mathrm{Tor}$ schéma des tores maximaux, $\mathrm{Kill}$ schéma des couples
de Killing
\item[$\mathrm{Bor}$] schéma des \struck{drapeaux} Borels,
\item[$\mathfrak{B}$] le « Borel universel » $\subset G_{\mathrm{Bor}}$, donc
$\mathfrak{B}/\mathfrak{B}_u\simeq T_{\mathrm{Bor}}$
\item[$\underline{\mathfrak{L}}$] son algèbre de Lie, qui est un fibré
vectoriel (\uncertain{et autres structures}) sur $\mathfrak{B}$\,,
$\underline{\mathfrak{L}}/\underline{\mathfrak{L}}_u\simeq
\underline{t}_{\mathrm{Bor}}$. On a un morphisme « surjectif »
$p\colon\underline{\mathfrak{L}}\to\underline{\mathfrak{g}}$ et un morphisme
lisse $\tilde f'\colon\underline{\mathfrak{L}}\to\underline{t}=\mathfrak{J}'$
\item[$\underline{\mathfrak{L}}^{\mathrm{rss}}$] sous-fibré des éléments
« rég.\ ss » dans $\underline{\mathfrak{g}}_{\mathrm{Bor}}$\,;
$=p^{-1}(\underline{\mathfrak{g}}^{\mathrm{rss}})$.
\end{itemize}
\note{« semi- » est ajouté au-dessus de « simples » aux deux endroits. Les
deux morphismes $p$ et $\tilde f'$ sont encadrés dans la marge droite,
d'une encre plus claire, et reliés par une accolade ; le premier $\sim$ de
$\tilde f'$ se lit mal. La ligne de $\mathrm{Tor}$ et $\mathrm{Kill}$ est
ajoutée à droite de $\mathfrak{J}^{\mathrm{rss}}$, liée par un trait à
« Kill ».}

On a un isom.\ canon.
\[
  \underline{\mathfrak{L}}^{\mathrm{rss}}\simeq \mathrm{Kill}\times_S
  \underline{t}^{\mathrm{rss}}
\]
et un carré cartésien

\begin{tikzcd}
\underline{\mathfrak{g}}^{\mathrm{rss}} \arrow[d] &
\underline{\mathfrak{L}}^{\mathrm{rss}}\simeq\mathrm{Kill}\times\underline{t}^{\mathrm{rs}} \arrow[l] \arrow[d] \\
\mathrm{Tor} & \mathrm{Kill} \arrow[l]
\end{tikzcd}

\marginal{$\operatorname{Lie}(\mathcal{T})^{\mathrm{rss}}$, où $\mathcal{T}$
est le tore maximal universel}
\note{la note marginale, à gauche, est reliée par une flèche au coin
$\underline{\mathfrak{g}}^{\mathrm{rss}}$ du carré.}

Il y a une opération naturelle de $W$ sur $\mathrm{Kill}$, faisant de
$\mathrm{Kill}$ un torseur sous $W_{\mathrm{Tor}}$, d'où une opération de $W$
sur le produit fibré $\underline{\mathfrak{L}}^{\mathrm{rss}}$, qui sur
$\mathrm{Kill}\times\underline{t}^{\mathrm{rs}}$ se réduit à l'action
\emph{diagonale}. Comme
$\underline{\mathfrak{g}}^{\mathrm{rss}}\simeq
\underline{\mathfrak{L}}^{\mathrm{rss}}/W$, on trouve un hom naturel
\[
  \underline{\mathfrak{g}}^{\mathrm{rss}}\longrightarrow
  \mathfrak{J}=\underline{t}^{\mathrm{rss}}/W
\]
\note{ce dernier alinéa est d'une encre brune, plus tardive que le reste de
la page.}

\page{5}
donnant naissance au diagramme

\begin{tikzcd}
\underline{\mathfrak{g}}^{\mathrm{rss}} \arrow[d, "f^{\mathrm{rss}}"'] &
\underline{\mathfrak{L}}^{\mathrm{rss}} \arrow[l, "q^{\mathrm{rss}}"'] \arrow[d, "\tilde f'^{\mathrm{rss}}"] \\
\mathfrak{J}^{\mathrm{rs}}=\underline{t}^{\mathrm{rss}}/W &
\mathfrak{J}'=\underline{t}^{\mathrm{rss}} \arrow[l]
\end{tikzcd}

Supposons que $W$ opère \struck{trivial.}\ \emph{librement} sur
$\underline{t}^{\mathrm{rss}}$\,; (ce qui est le cas \uncertain{si} les car.\
résiduelles sont toutes $\neq 2$.
\marginal{Il revient ici au même qu'il opère fidèlement \struck{\ill{}} sur
les fibres géom.\ de $\underline{t}$,}
Alors le diagramme précédent est cartésien\,; \struck{\ill{}}
\note{la note, dans une bulle en haut à droite, est reliée à
« librement » ; la parenthèse ouverte avant « ce qui » ne se ferme pas.}

Je dis que le morphisme
$\underline{\mathfrak{g}}^{\mathrm{rss}}\to\mathfrak{J}^{\mathrm{rss}}$ se
prolonge en un morphisme $\underline{\mathfrak{g}}\to\mathfrak{J}$
(évidemment unique), d'où un diagramme commutatif \ill{}

\begin{tikzcd}
\underline{\mathfrak{g}} \arrow[d, "f"'] & & \underline{\mathfrak{L}} \arrow[ll, "q"'] \arrow[dl, "\tilde f'"] \\
\mathfrak{J}=\underline{t}/W & \mathfrak{J}'=\underline{t} \arrow[l] &
\end{tikzcd}

\note{sous l'avant-dernier $\underline{t}$ un exposant « rss » est biffé.}

\marginal{NB L'anneau affine de $\mathfrak{J}$ sur $S$ est le sous-anneau des
invariants de l'anneau \ill{} de $\underline{\mathfrak{g}}$ sur $S$, et $f$
est \ill{} \ill{} d'un \ill{} dans \ill{} invariants dans un schéma affine
(\uncertain{construction probablement} \ill{} \ill{} sur un schéma
\ill{}~?)}
\note{bloc encadré dans la marge gauche, d'une écriture serrée et en partie
effacée ; seul le début se lit sûrement.}

Dém.\ de l'existence de $f$. On se ramène au cas $S=\operatorname{Spec}
\mathbf{Z}$, $G$ déployé. \struck{Alors il suffit de \ill{} pour \ill{} \ill{}
de \ill{}.} Comme $\underline{\mathfrak{g}}$ est normal, $\mathfrak{J}$
affine, il suffit de le voir sur les fibres géométriques\,; là c'est connu
(Kostant\,…).

Introduisons
\[
  \underline{\mathfrak{g}}'=\underline{\mathfrak{g}}\times_{\mathfrak{J}}\mathfrak{J}'
\]
on trouve donc un diagramme

\begin{tikzcd}
\underline{\mathfrak{g}} \arrow[d, "f"'] & \underline{\mathfrak{g}}' \arrow[l] \arrow[d, "f'"'] & \underline{\mathfrak{L}} \arrow[l, "p"'] \arrow[dl, "\tilde f'"] \arrow[r] \arrow[ll, bend right=30, "q"'] & D \\
\mathfrak{J} & \mathfrak{J}' \arrow[l] & &
\end{tikzcd}

\note{la flèche $q$ est tracée en arc au-dessus de la ligne. Le but $D$ de la
dernière flèche, écrit à l'encre pâle, est de lecture douteuse : sans doute
le schéma de base de $\underline{\mathfrak{L}}$. Deux ovales au crayon pâle
entourent $\underline{\mathfrak{g}}'$ et $\underline{\mathfrak{L}}$.}

avec $p$ morphisme projectif, $\tilde f'$ morphisme lisse\,; dans le cas
favorable où $W$ opère \struck{fidèlement} \emph{librement} sur
$\mathfrak{J}^{\mathrm{rss}}$, i.e.\ fidèlement sur les fibres géom.\ de
$\underline{t}$, alors $p$ \struck{induit} est un isom.\ birationnel, plus
précisément il induit un isom.\
$\underline{\mathfrak{L}}^{\mathrm{rss}}\simeq
\underline{\mathfrak{g}}'^{\mathrm{rss}}$, où
$\underline{\mathfrak{g}}'^{\mathrm{rss}}=\underline{\mathfrak{g}}'\,|\,
\underline{\mathfrak{g}}^{\mathrm{rss}}$.
\note{ce dernier alinéa, d'encre brune, est plus tardif.}

\page{6}
Par réduction au cas de la base $\mathbf{Z}$ ou $\mathbf{Z}[1/2]$, et par
application du Main theorem, on trouve que le plus grand ouvert de
$\underline{\mathfrak{g}}'$ sur lequel $p$ est un isom.\ est formé des
$x'\in\underline{\mathfrak{g}}'$ tels que $p^{-1}(x')$ soit fini (ou encore,
\add{contienne un} élément isolé) i.e.\ l'image inverse de l'\uncertain{ouvert}
$\underline{\mathfrak{g}}^{\mathrm{reg}}\subset\underline{\mathfrak{g}}$,
formé des points $x$ tels que $q^{-1}(x)$ soit fini (ou encore\,…), i.e.\
\uncertain{les} $x\in\underline{\mathfrak{g}}$ contenus dans un nb
\emph{fini} seulement de sous-algèbres de Borel. En vertu de Springer PM
n° 30, 5.3, \struck{il résulte} \add{on voit} que cet
ouvert \struck{contient} \uncertain{rencontre} les points unipotents de toute
fibre, et on prévoit que l'on \uncertain{pourra} en déduire qu'il rencontre
toute fibre de $f\colon\underline{\mathfrak{g}}\to\mathfrak{J}$\,?
\textbf{Non} en car.~2, cas $A_1$, $B_2$\,… peut-être d'autres\,…
\note{« Non » est souligné deux fois et écrit plus tard, à l'encre plus
noire, comme le $\mathfrak{J}$ repassé qui le précède.}

\textbf{Questions}.
\begin{enumerate}
\item[0)] $W$ opère-t-il fidèlement sur $\underline{t}^{\mathrm{rss}}$\,?
\item[a)] $\mathfrak{J}$ est-il régulier, i.e.\ $\mathfrak{J}'$ plat sur
$\mathfrak{J}$\,?
\item[b)] $f$ est-il plat (indépendamment de a))\,; en fait
$f$ plat $\Longleftrightarrow$ $\mathfrak{J}$ régulier et $f$
équidimensionnel.
\marginal{Oui pour les « bons » nombres premiers\,?}
\item[c)] Les fibres de $f$ géom.\ \struck{\ill{}} \uncertain{intègres}\,?
Sont-elles des \ill{}, \uncertain{la} \ill{} l'adhérence d'une orbite sous
$G$, \uncertain{le} complémentaire \ill{} de codim.\ $\geqslant 2$ (car
d'ailleurs dans la fibre \ill{} de \ill{} normales, donc
\uncertain{équivaut} que les fibres de $f$ \uncertain{soient normales}
(\uncertain{si bons nombres}), \ill{} $f'$ normal, donc
$\underline{\mathfrak{g}}'$ normal, dans
$\underline{\mathfrak{L}}\to\underline{\mathfrak{g}}'\to\underline{\mathfrak{g}}$
factorisation de Stein)
\item[d)] Les points de lissité de $f$ sont-ils les points réguliers
(\uncertain{quasi-réguliers}\,?) i.e.\ ceux au dessus desquels $\mathfrak{J}$
\ill{}\,? \struck{\ill{} \ill{}} Sont-ce aussi les \ill{} \ill{} dans
\ill{} de dimension \ill{} \ill{} \ill{} minimum, soit $r$\,?
\item[e)] \struck{Les orbites de $G$ sur $\underline{\mathfrak{g}}$ et
$\mathfrak{J}$} Définir stratification orbitale \ill{} \ill{} Chevalley sur
$\mathbf{Z}$, suivant une filtration naturelle de $\underline{t}$\,; cette
structure orbitale devrait être homologue (\uncertain{$G$ opérant sur
$\tilde G$}) à la structure orbitale de $G$ \uncertain{opérant sur}
$\tilde G$ \ill{} connue)\,; relations entre les \uncertain{classes} \ill{}
et orbites combinatoires correspondantes\,: des orbites \uncertain{unipotentes}\,?
\,/ Questions analogues pour $\mathfrak{B}$, $\underline{\mathfrak{L}}$.
\end{enumerate}
\marginal{f) Reprendre théorie de Kostant\,: A-t-on section de l'ouvert des
pts réguliers etc\,…}
\marginal{g) Étudier fibres de $p\colon\underline{\mathfrak{L}}\to
\underline{\mathfrak{g}}'$, i.e.\ les composantes connexes de la variété des
sous-alg.\ de Borel de $\underline{\mathfrak{g}}$ contenant un élément donné.
Quelle est sa dimension (relation avec dim.\ de l'orbite)\,?}
\note{la page est très chargée et d'écriture rapide ; les questions c), d)
et e) ne se lisent que par fragments. Les items f) et g) sont écrits dans la
marge gauche, en face de e), à l'encre brune ; e) est lui-même en partie à
l'encre brune. Dans la marge gauche, en face de c), un bloc commençant par
« NB » et contenant « $\underline{t}/W$ » et
« $\underline{\mathfrak{g}}$ sur $S$ » est entièrement biffé de hachures et
ne se lit pas.}

\section*{Lissification des morphismes de Steinberg}

\note{titre de sa main, en haut à droite du feuillet 7, qui ne porte rien
d'autre : une chemise, sans \texttt{page}.}

\page{8}
\begin{itemize}
\item[$G$] schéma en groupes réductif sur base $S$.
\item[$D=\mathrm{Bor}(G)$,] schéma des Borels.
\item[$\mathfrak{B}\to D$] le « Borel universel » ; $\mathfrak{B}\subset
G_D=G\times_S D$ sous-schéma en groupes
\item[$\mathfrak{B}_u\subset\mathfrak{B}$] la partie unipotente
\item[$\mathfrak{B}\to T_D=\mathfrak{B}/\mathfrak{B}_u$] tore sur $D$.
\end{itemize}

On sait que $T_D$ est muni d'une donnée de descente de $D$ à $S$, d'où un tore
$T$ sur $S$ \struck{\ill{}} avec isom.\
\[
  \struck{\ill{}}\ T_D\simeq T\times_S D, \quad\text{donc}\quad
  \hat T_D\simeq M_D,\ \text{où}\ M=\hat T\ \text{réseau sur}\ S.
\]
\begin{itemize}
\item[$K=\mathrm{Kill}(G)$,] schéma des couples de Killing $(T,\mathfrak{B})$,
$K=\underline{\mathrm{Tmax}}(\mathfrak{B})$
\item[$K\to D$] \uncertain{le} morphisme standard
\item[$T_K\hookrightarrow\mathfrak{B}_K$] le \struck{tore maximal} couple de
Killing universel ; $T_K=T\times_S K$
\item[$N_K=\operatorname{Norm}(T_K,G_K)$]
\item[\struck{$P_K=\ill{}$}]
\item[$W_K=N_K/T_K$] opère sur $T_K$, $M_K$
\end{itemize}
\note{à gauche de « $=\mathrm{Kill}(G)$ » une lettre est biffée et
remplacée par $K$. Le couple $(T,\mathfrak{B})$ est ajouté au-dessus de
« Killing », avec un trait vers $K=\mathrm{Tmax}(\mathfrak{B})$ ; un trait
vertical, dans la marge droite, tient les lignes de $K\to D$ à $N_K$.}

On sait que ces opérations proviennent d'opérations d'un groupe $W/S$ sur
$T$, $M$ (avec sa structure \ill{})
\[
  W_K\simeq W\times_S K
\]
\struck{\ill{}} \add{soit} $I=T/W$.

On définit un \uncertain{morphisme} \add{unique} (canonique), invariant par
automorphismes intérieurs (donc par tous automorphismes de $G$, quand on fait
opérer sur $I$ les automorphismes extérieurs)
\[
  G\xrightarrow{\ f\ } I
\]
La restriction de $f$ à $T_0$, si $T_0\subset B_0$, est donnée par
\[
  f\,|\,T_0 = \text{morph.\ can.}\ T_0\to T_0/W_0\simeq I .
\]
\marginal{Défini sans hypothèse $G$ simpl.\ connexe\,? \emph{Oui}
\uncertain{en dehors de} \ill{} (\uncertain{Steinb.}) $f$ normal en tous
cas\,?? [Mais alors il faudrait que $I$ soit lisse$/S$ en tous cas\,!]}
\note{la note, à l'encre brune dans la marge gauche, est reliée par une
accolade à « $G\xrightarrow{f}I$ » ; « Un », entouré, est écrit devant « La
restriction ».}

\page{9}
Soit $G'$ défini par la condition d'avoir un diagramme cartésien

\begin{tikzcd}
\mathfrak{X}=G \arrow[d, "f"'] & \mathfrak{X}'=G' \arrow[l] \arrow[d, "f'"] & \mathfrak{B}=\tilde{\mathfrak{X}} \arrow[l, "p"'] \arrow[dl, "\tilde f"] \\
I & I'=T \arrow[l, "\pi"] &
\end{tikzcd}

\note{« cart » est écrit, souligné, au centre du carré de gauche ;
$\mathfrak{X}'=G'$ et $\mathfrak{B}=\tilde{\mathfrak{X}}$ sont encadrés
ensemble.}

Utilisant le fait que $f$ est normal, donc $f'$ normal, on trouve que $G'$
est normal.

On va définir un hom
\[
  p\colon \mathfrak{B}\longrightarrow G'
\]
i.e.\ des hom
\[
  \tilde f\colon\mathfrak{B}\to T,\qquad h\colon\mathfrak{B}\to G
  \qquad\text{tels que }\ \pi\tilde f=f h
\]
On prendra
\[
  \begin{cases}
  h(g,B)=g\\
  \tilde f(g,B)=g\in B/B_u(-)\simeq T(-),
  \end{cases}
\]
i.e.\ $h$ est le composé
$\mathfrak{B}\to G_D=G\times_S D\xrightarrow{\mathrm{pr}_1}G$, i.e.\ $\tilde f$
est le composé
$\mathfrak{B}\to\mathfrak{B}/\mathfrak{B}_u=T_D\simeq T\times_S
D\xrightarrow{\mathrm{pr}_1}T$.
\note{les deux « i.e. » sont écrits en face des deux lignes de l'accolade ;
ils sont reportés ici sous elle.}

Pour voir que $\pi\tilde f=fh$, il suffit de le voir sur l'ouvert dense
\add{$\mathfrak{B}^{\mathrm{rs}}$} des \struck{pa} $(g,B)$ avec $g$ régulier
ss, \struck{\ill{}} donc $g\in T\subset B$, et alors c'est trivial.
\note{« $\mathfrak{B}^{\mathrm{rs}}$ des » est écrit au-dessus de la ligne,
à la place d'un mot biffé.}

\textbf{NB}
$\begin{cases}\text{$h$ est propre, donc $p$ est propre (comme on verra qu'il est birationnel)}\\ \text{$\tilde f$ est lisse.}\end{cases}$

\begin{itemize}
\item[$G^{\mathrm{reg}}$] l'ouvert des points réguliers de $G$ $=$ ouvert
des pts de lissité de $f$
\item[$G'^{\mathrm{reg}}$] son image inverse dans $G'$
\item[$\mathfrak{B}^{\mathrm{reg}}$] son image inverse dans $\mathfrak{B}$.
\end{itemize}
\note{« $=$ ouvert des pts de lissité de $f$ » est ajouté en fin de ligne,
d'une écriture plus petite.}

\page{10}
\textbf{Théorème}
\begin{enumerate}
\item[(i)] (Pour mémoire) $p\colon\mathfrak{B}\to G'$ est projectif,
$\tilde f\colon\mathfrak{B}\to T$ est lisse.
\item[(ii)] $p$ induit
$\mathfrak{B}^{\mathrm{reg}}\overset{\mathrm{df}}{=}\mathfrak{B}\,|\,G'^{\mathrm{reg}}
\xrightarrow{\ \sim\ }G'^{\mathrm{reg}}$
et $G'^{\mathrm{reg}}$ est le plus grand ouvert au dessus duquel $p$ est un
isom.
\item[(iii)] (Pour mémoire) $\mathfrak{B}-\mathfrak{B}^{\mathrm{reg}}$ est de
codim.\ $\geqslant 2$ \struck{\ill{}} dans $\mathfrak{B}$ (fibre par fibre sur
$S$ \uncertain{ou} sur $D$)\,; $G'-G'^{\mathrm{reg}}$ a toutes ses
composantes irréductibles de codimension $3$ (fibre par fibre sur $S$)
\item[(iv)] (Si $S$ \struck{\ill{}} normal \uncertain{loc.\ noeth.})
\[
  \operatorname{Pic}(\mathfrak{B})\xrightarrow{\ \sim\ }
  \operatorname{Pic}(\mathfrak{B}^{\mathrm{reg}})\xleftarrow{\ \sim\ }
  \operatorname{Pic}(G'^{\mathrm{reg}})
  \qquad
  \operatorname{Pic}(D)\xrightarrow{\ \wr\ }\operatorname{Pic}(\mathfrak{B})
\]
\marginal{[\struck{\ill{}} Si $T$ \uncertain{est trivial}, \struck{ou} \ill{}
\uncertain{génériquement} si \uncertain{\ill{} les points maximaux de $S$,}
l'image de $\pi_1(S,\xi)$ dans $\operatorname{Ext}(G_\xi)$ est telle que
$H^1(\Gamma,M)=0$] \uncertain{on a une réciproque} si $S$ régulier.}
et en faisceautisant sur $S$ (pour la top.\ étale, disons)
\[
  \underline{\mathrm{Pic}}_{\mathfrak{B}/S}\xrightarrow{\ \sim\ }
  \underline{\mathrm{Pic}}_{\mathfrak{B}^{\mathrm{reg}}/S}\xleftarrow{\ \sim\ }
  \underline{\mathrm{Pic}}_{G'^{\mathrm{reg}}/S}
  \qquad
  M\to P\simeq\underline{\mathrm{Pic}}_{D/S}\xrightarrow{\ \wr\ }
  \underline{\mathrm{Pic}}_{\mathfrak{B}/S}
\]
\note{sur la page, $\operatorname{Pic}(D)$ et $M\to P\simeq
\underline{\mathrm{Pic}}_{D/S}$ sont écrits sous le premier terme, avec une
flèche verticale vers le haut marquée $\wr$ ; ils sont ramenés ici sur la même
ligne. Sous $M\to P$ : « (injectif si $G$ semi-simple) ». Une fin de ligne
après $\operatorname{Pic}(G'^{\mathrm{reg}})$ et un terme sous
$\underline{\mathrm{Pic}}_{\mathfrak{B}/S}$ sont biffés de hachures. La
parenthèse de condition, sous la flèche de gauche, est une insertion de
lecture difficile.}
De plus, le cône ample pour $\underline{\mathrm{Pic}}_{\mathfrak{B}/S}$ est
le \uncertain{cône} \ill{} pour $\underline{\mathrm{Pic}}_{D/S}$, qui est le
cône des poids $P^{+}$. \struck{\ill{}}
\end{enumerate}
\note{un trait vertical, dans la marge gauche, tient tout l'énoncé du
théorème.}

\textbf{Démonstration} (i) et (iii) connus (pour (iii) réf.\,: Steinberg\,…)

(ii) On commence à voir, posant \struck{$G^{\mathrm{rs}}=\ill{}$}
\[
  \begin{aligned}
  &G^{\mathrm{rs}}\subset G^{\mathrm{reg}} && \text{ouvert des pts de $G$ qui
  sont réguliers semi-simples}\\
  &G'^{\mathrm{rs}}\subset G'^{\mathrm{reg}},\ \
  \mathfrak{B}^{\mathrm{rs}}\subset\mathfrak{B}^{\mathrm{reg}} && \text{images
  inverses,}
  \end{aligned}
\]
que
\[
  \mathfrak{B}^{\mathrm{rs}}\xrightarrow{\ \sim\ }G'^{\mathrm{rs}}
\]
D'autre part, on sait par Steinberg que les fibres de
$\mathfrak{B}^{\mathrm{reg}}\to G'^{\mathrm{reg}}$ sont finies, d'ailleurs c'est
un morphisme propre. Comme $G'^{\mathrm{reg}}$ est normal, on conclut par le
Main Theorem.
\marginal{On sait aussi par Steinberg que les fibres des \uncertain{centralisateurs}
\uncertain{plats} \ill{} \ill{} \ill{}.}
\note{les trois dernières lignes, et la note qui les prolonge en bas à
droite, sont à l'encre brune.}

\page{11}
(iv) Supposons \struck{d'abord} $S$ \struck{régulier} \add{normal},
\struck{Alors $\operatorname{Pic}(\mathfrak{B})\simeq
\operatorname{Pic}(\mathfrak{B}^{\mathrm{reg}})$ résulte}
$\operatorname{Pic}(\mathfrak{B})\to
\operatorname{Pic}(\mathfrak{B}^{\mathrm{reg}})$ injectif résulte des
considérations de \uncertain{profondeur} et \add{donc aussi tous les schémas
écrits.} \struck{\ill{}} du fait que
$\mathfrak{B}-\mathfrak{B}^{\mathrm{reg}}$ est de codim.\ $\geqslant 2$,
\struck{\ill{}} fibre par fibre. Le fait que
$\operatorname{Pic}(D)\to\operatorname{Pic}(\mathfrak{B})$ \add{et
$\operatorname{Pic}(D)\to\operatorname{Pic}(\mathfrak{B}^{\mathrm{reg}})$}
\struck{\ill{}} \add{surjectifs} résulte de EGA IV 21.4.13 (EGA\,IV 53),
compte tenu du

\textbf{Lemme} $G$ groupe \struck{\ill{}} \add{semi-simple}
simplement \uncertain{connexe} \add{sur un corps $k$},
$B$ un Borel, alors
$\operatorname{Pic}(B)$ \struck{\ill{}}
$=\operatorname{Pic}(B^{\mathrm{reg}})=\operatorname{Pic}(T)=H^1(\Gamma,M)$
$(T=B/B_u)$, $M=\hat T$)
\marginal{On a $\operatorname{Pic}(B)\simeq\operatorname{Pic}(T)$}
\note{le lemme est tenu par un double trait vertical à gauche ; ses termes
sont ajoutés en plusieurs couches au-dessus de la ligne. La note marginale
est à gauche, en face de la preuve.}

En tout cas, utilisant \struck{la relation} $B\to T=B/B_u$, on se ramène au
cas où $\operatorname{Pic}(T)=0$ (NB $B_u$ a
\uncertain{dévissage} en groupes additifs). Or on sait que
$T=\prod_{k'/k}T'$, $T'$ un tore sur une extension \uncertain{galoisienne}
finie \uncertain{déployée}, \add{galoisienne (groupe $\Gamma$)}
\struck{Donc on a \ill{} équiv.\ de cat.}
\[
  \struck{\underline{\mathrm{Iso}}(T)\simeq\underline{\mathrm{Iso}}(T')}
  \ \text{donc}
\]
\[
  0\to H^1(\Gamma,\underbrace{H^0(T',\underline{O}^{*}_{T'})}_{M\times k'^{*}})
  \to\operatorname{Pic}(T)\to
  H^0(\Gamma,\underbrace{H^1(T',\underline{O}^{*}_{T'})}_{=0})
\]
\note{à gauche de la suite, un $\operatorname{Pic}(T)$ biffé ; à droite du
premier terme, une glose biffée de hachures.}
donc on trouve
\[
  \operatorname{Pic}(T)=H^1(\Gamma,M)\times
  \underbrace{H^1(\Gamma,k'^{*})}_{=0\ \text{H.~90}}
  =H^1(\Gamma,M)=H^1(k,\underline{M})
\]
Si $P$ est tordu, cela peut être $\neq 0$\,; si $B$ est la fibre générique de
$\mathfrak{B}$ sur $D$ ($S$ \uncertain{réduit}), alors cela indique donc que
l'hom\ $\operatorname{Pic}(D)\to\operatorname{Pic}(\mathfrak{B})$ n'est pas
surjectif. L'injectivité est claire, car $\mathfrak{B}$ a une section sur $D$
(la section unité). En faisceautisant pour la topologie étale, on trouve donc
le diagramme d'isom pour les $\underline{\mathrm{Pic}}$. / Comme
$\mathfrak{B}\to D$ est affine, et a une section, \ill{} et
\uncertain{l'on voit} qu'un faisceau inv.\ sur $D$ est ample ssi son image
inverse sur $\mathfrak{B}$ l'est\,:
\marginal{indép.\ de l'hypothèse normal sur $S$\,!}

\struck{Peut-on se passer dans le cas non régulier\,? a) L'isomorphisme
\ill{} $\operatorname{Pic}(D)\simeq\operatorname{Pic}(\mathfrak{B})$
\uncertain{et} \ill{} \ill{} $\operatorname{Pic}(D)\to
\operatorname{Pic}(\mathfrak{B}^{\mathrm{reg}})$ \uncertain{dès que} $S$ est
normal. L'injectivité \ill{} \ill{} évidemment\,!}
\note{bloc encadré et barré de traits obliques, en bas de page.}

\page{12}
\textbf{Remarques} a) Si $T$ peut-être tordu, on ne peut plus affirmer que
l'hom injectif
\[
  \operatorname{Pic}(D)\longrightarrow\operatorname{Pic}(\mathfrak{B})
\]
soit surjectif. Mais (si $S$ \struck{\ill{}} n'a qu'un nb fini de
composantes connexes) on voit que le \emph{conoyau est fini}\,; en fait
\ill{} \ill{} pour $S$ irréductible \uncertain{et} \uncertain{normal}
\ill{} \uncertain{complète}
\[
  0\to\operatorname{Pic}(D)\to\operatorname{Pic}(\mathfrak{B})\to
  H^1(\Gamma,M_{\bar\xi}),
  \qquad
  \operatorname{Pic}(\mathfrak{B}_\eta)\simeq\operatorname{Pic}(T_\eta)
  \simeq\operatorname{Pic}(T_\xi)
\]
où $\eta$ pt gén.\ de $D$, $\xi$ pt gén.\ de $S$, où $\Gamma$ est l'image de
$\pi_1(S,\bar\xi)$ dans $\operatorname{Aut}(M_{\bar\xi})$. Si $S$ est
régulier, \struck{\ill{}} on peut mettre une \uncertain{flèche} à droite,
i.e.\ la condition $H^1(\Gamma,M_{\bar\xi})=0$ est néc.\ et suff.\ pour que
l'on ait \struck{\ill{}} \add{on a aussi} surjectivité
$\operatorname{Pic}(D)\to\operatorname{Pic}(\mathfrak{B})$\,; \struck{\ill{}}
surjectivité $\operatorname{Pic}(D)\to
\operatorname{Pic}(\mathfrak{B}^{\mathrm{reg}})$
[car $\operatorname{Pic}(\mathfrak{B}_\xi)\to
\operatorname{Pic}(\mathfrak{B}_\xi^{\mathrm{reg}})$ est \uncertain{injectif}
\add{en tous cas} \struck{surjectif} par des raisons de codim.\ $\geqslant 2$
et de régularité]
\note{la ligne des isomorphismes $\operatorname{Pic}(\mathfrak{B}_\eta)$ …
est écrite sous la flèche, avec une glose « où $\eta$ pt gén.\ de $D$, $\xi$
pt gén.\ de $S$ » ; elles sont réunies ici.}

b) L'injectivité de $\operatorname{Pic}(D)\to\operatorname{Pic}(\mathfrak{B})$
et la caractérisation des éléments de $\operatorname{Pic}(D)$ amples$/S$
comme ceux dont l'image inverse sur $\mathfrak{B}$ est ample$/S$,
\struck{\ill{}} \uncertain{restent} vrais sans hyp.\ de normalité sur $S$.
\add{et $H^1(\xi,M)=0$ (p.\ ex.\ $T$ non tordu)}

c) Supposons $S$ normal. Alors
\[
  \operatorname{Pic}(S)\xrightarrow{\ \sim\ }\operatorname{Pic}(G')
\]
En effet, on \struck{sait que $\operatorname{Pic}(G')$} a
\uncertain{injectivité} car $G'$ a une section sur $S$ (section unité). Reste
à prouver que si $\xi\in\operatorname{Pic}(G')$ induit $0$ sur la section
unité, alors il est nul. Or
\[
  \operatorname{Pic}(G')\hookrightarrow
  \operatorname{Pic}(G'^{\mathrm{reg}})\,[\simeq
  \operatorname{Pic}(\mathfrak{B}^{\mathrm{reg}})\leftleftarrows
  \operatorname{Pic}(\mathfrak{B})]
\]
\uncertain{injectif} par raisons de profondeur, donc il suffit de voir que
l'image \struck{inverse} de $\xi$ dans $\operatorname{Pic}(\mathfrak{B})$ est
nulle. Mais il suffit de voir que la restriction de cette dernière à la
section unité de $\mathfrak{B}$ sur $D$ est nulle. Mais $D\to G'$
\note{la phrase se poursuit au feuillet suivant. Le symbole
$\xi$ de ce paragraphe (une accolade ouvrante sur la page) est de lecture
douteuse et n'est pas le point générique de a). Les trois dernières lignes
sont à l'encre brune.}

\page{13}
se factorise par la section unité de $G'$ sur $S$\,! En fait, sans hyp.\ sur
$T$, on obtient des suites exactes

\begin{tikzcd}
G'^{\mathrm{reg}} \arrow[d, hook] & \mathfrak{B}^{\mathrm{reg}} \arrow[l] \arrow[d, hook] \\
G' \arrow[d] & \mathfrak{B} \arrow[l] \\
S \arrow[u, bend left=40] & D \arrow[u, hook] \arrow[ul] \arrow[l]
\end{tikzcd}

\[
  \cdots\to\operatorname{Pic}(S)\to\operatorname{Pic}(G')\longrightarrow
  \underbrace{H^1(\xi,M)}_{\text{groupe fini}}
\]
\note{le diagramme est dessiné dans la marge gauche, la suite à droite ; la
flèche $D\to\mathfrak{B}$ est la section unité, et la flèche courbe de $S$
vers $G'$ la section unité de $G'$.}

d) On a \struck{$\operatorname{Pic}(G)$} ($S$ normal \uncertain{irréd.}
\uncertain{de point générique} $\xi$)
\[
  0\to\operatorname{Pic}(S)\to\operatorname{Pic}(G)\to\operatorname{Pic}(G_\xi)
\]
\note{sous la première flèche : « injectif car $\exists$ section unité » ; sous
le dernier terme : « groupe \struck{fini tors.} qui est nul si $G_\xi$ est
simplement connexe et déployé ».}

\struck{\ill{}} D'autre part \struck{si $S$ régulier} on a
$\operatorname{Pic}(G)\xrightarrow{\ \text{inj}\ }
\operatorname{Pic}(G^{\mathrm{reg}})$ \add{sans hyp.\ de normalité,} (car on
a codim.\ $\geqslant 3$ fibre par fibre \add{sur $S$} \struck{de}
$G-G^{\mathrm{reg}}$), donc on trouve
\[
  \operatorname{Coker}\bigl(\operatorname{Pic}(S)\to
  \operatorname{Pic}(G^{\mathrm{reg}})\bigr)\ \text{est fini (et $=$ nul si
  $G_\xi$ est simplement connexe et déployé)}
\]
Cela montre \struck{bien} qu'on \uncertain{ne peut} trouver une résolution $\tilde X$ des
singularités de $f$ analogue à celle de $f'$, sans toucher à
$G^{\mathrm{reg}}$ (l'image inverse de $G-G^{\mathrm{reg}}$ dans $\tilde X$
étant de codim.\ $\geqslant 2$ fibre par fibre\,…)
\note{un « e) » isolé, biffé, figure en bas de page ; il n'est suivi de rien.}

\section*{Résumé th.\ de conjugaison Kostant}

\note{titre de sa main, en haut à droite du feuillet 14, chemise de carton
qui ne porte rien d'autre ; il se lit « Résumé th.\ de conjugaison Kostant
+ lettre à Kostant ». La lettre n'est pas dans ce lot. Les pages 15 à 19 sont
un tapuscrit en français, frappé à la machine, que sa main corrige et
complète : les lettres gothiques, que la machine ne portait pas, sont
ajoutées à l'encre ($\mathfrak{g}$ pour l'algèbre de Lie ; $\underline{G}_m$,
$\underline{\varnothing}$ sont soulignés à la frappe). Les fautes de frappe
recouvertes de \emph{x} ne sont pas relevées.}

\page{15}
Résumé de quelques résultats de Kostant (sous groupes simples de rang 1).
\marginal{(The principal 3-dimensional subgroup …, Am.\ Jour.\ Oct.\ 1959,
3.4, 3.5, 3.6, 4.2, 9.3)}

$G$ est un groupe semi-simple sur un corps $k$ de car.\ nulle,
$\mathfrak{g}$ son algèbre de Lie, \add{avec $e\neq 0$, $e$ nilpotent.}

\textbf{Théorème} 1. (i) Soient $x,e\in\mathfrak{g}$\,; pour tout entier $k$,
soit $\mathfrak{g}(k)$ le sous-espace propre de \ill{} correspondant à la
valeur propre $k$. Donc si $x=\operatorname{Lie}(i)(1)$, où
$i\colon\underline{G}_m\to G$, les $\mathfrak{g}(k)$ sont les composants de
$\mathfrak{g}$ suivant les caractères de $\underline{G}_m$, opérant sur
$\mathfrak{g}$ via $i$ et la représentation adjointe. Les conditions
suivantes sont équivalentes\,:
\note{« avec $e\neq 0$, $e$ nilpotent » est ajouté à la main sous la
première ligne et relié par une accolade à « Soient $x,e$ ». Après
« propre de » la frappe laisse un blanc que la main n'a pas rempli : on
attend $\operatorname{ad}(x)$.}

\begin{enumerate}
\item[a)] $e\in\mathfrak{g}(2)$ i.e.\ $[x,e]=2e$, et \add{pour tout entier
$k\geqslant 1$,} $\operatorname{ad}(e)^{k}\colon\mathfrak{g}(-k)\to
\mathfrak{g}(k)$ est un isom.
\item[b)] $e\in\mathfrak{g}(2)$, $\operatorname{ad}(e)\colon\mathfrak{g}(0)
\to\mathfrak{g}(-2)$ est surjectif, \add{et $x$ est de la forme $f(x_0)$, pour
$f\colon\mathrm{Sl}(2,k)\to G$ (\emph{élément de Dynkin}).}
\item[\struck{b')}] \struck{Comme b), avec
$\operatorname{ad}(e)\colon\mathfrak{g}(0)\to\mathfrak{g}(2)$ \ill{}}
\item[c)] $e\in\mathfrak{g}(2)$, et $x\in\operatorname{ad}(e)(\mathfrak{g})$.
\item[c')] $e\in\mathfrak{g}(2)$, et $x\in\operatorname{ad}(e)(\mathfrak{g}(-2))$.
\item[d)] Il existe un homomorphisme $f\colon\mathrm{Sl}(2,k)\to G$ tel que
$f'(e_0)=e$, $f'(x_0)=x$
\end{enumerate}
où $(x_0,e_0,f_0)$ est la base de $\operatorname{Lie}\mathrm{Sl}(2,k)$
correspondant à l'épinglage habituel de ce groupe.
\note{la ligne b') est entièrement barrée à la main ; le complément de b) est
écrit à sa suite, et relié à b) par un trait. Dans d), les accents de $f'$
sont ajoutés à la main.}

Ces conditions impliquent qu'il existe un homomorphisme
$i\colon\underline{G}_m\to G$, tel que $\operatorname{Lie}(i)(1)=x$, a
fortiori que $x$ est semi-simple. De plus, le $f$ de d) est déterminé de façon
unique.

(ii) Pour tout élément nilpotent $e$ de $\mathfrak{g}$, il existe un
homomorphisme $f\colon\mathrm{Sl}(2,k)\to G$ tel que $f(e_0)=e$, ou ce qui
revient au même, tel que $f(u_0)=u$ (où $u=\exp(e)$, $u_0=\exp(e_0)$)\,;
\add{$f$ est uniquement déterminé par $x=f(x_0)$ i.e.\ par $i=f\circ i_0$.}
En vertu de (i), il revient au même, lorsque $e\neq 0$ i.e.\ $u\neq 1$, de
dire qu'il existe un $x\in\mathfrak{g}$ tel que $(x,e)$ satisfasse aux
conditions équivalentes de (i), ou encore un homomorphisme
$i\colon\underline{G}_m\to G$ tel que $(\operatorname{Lie}(i)(1),e)$ y
satisfasse. Les homomorphismes $f$ en question sont conjugués par
$Z=\operatorname{Cent}(e)=\operatorname{Cent}(u)$ (i.e.\ pour $e\neq 0$, les
$x$ en question, ou encore les $i$ en question, sont conjugués par
$\operatorname{Cent}(u)$) et lorsque $e\neq 0$, les $f$ en question (ou encore
les

\page{16}
$x$ en question, ou les $i$ en question) forment un \emph{torseur} sous le
groupe $Z^{+}=Z\cap U$ \struck{$\operatorname{Cent}(u)$}, qui est un groupe
unipotent\struck{ (en fait, celui introduit dans (iv))}. Enfin, supposant
toujours $e\neq 0$, l'ensemble des $x$ en question est un torseur sous
\struck{$\mathfrak{g}^{e}=$ (centralisateur de $e$ dans $\mathfrak{g}$)
$=\operatorname{Ker}\operatorname{ad}(e)=\operatorname{Lie}
\operatorname{Cent}(u)$}
$\operatorname{Lie}(Z^{+})=\mathfrak{g}^{e}\cap\operatorname{Fil}^{1}
(\mathfrak{g})$ (\add{torseur} pour l'addition, i.e.\ est de la forme
$a+\mathfrak{g}^{e}$, pour un $a\in\mathfrak{g}$ convenable).
\struck{et on a $\mathfrak{g}^{e}\subset\sum_{k\geqslant 1}\mathfrak{g}(k)
=\operatorname{Fil}^{1}(\mathfrak{g})$}
\marginal{cf.\ plus bas définition de $U$. On a $Z\simeq Z_0\cdot Z^{+}$
\struck{\ill{}} (produit semi-direct) avec
$Z_0=Z\cap L=\operatorname{Cent}(f)$.}
\note{les ajouts à la main de ce passage sont d'une même encre : $Z^{+}=$
$Z\cap U$ remplace un « $\operatorname{Cent}(u)$ » dactylographié et biffé,
$\operatorname{Lie}(Z^{+})=\mathfrak{g}^{e}\cap\operatorname{Fil}^{1}$
remplace la définition dactylographiée de $\mathfrak{g}^{e}$, biffée. La
note marginale est écrite en oblique, à gauche.}

(iii) Soit $x$ \add{$\neq 0$} un élément de Dynkin de $\mathfrak{g}$, donc
associé à un homomorphisme de Dynkin $i\colon\underline{G}_m\to G$. Alors les
éléments \add{$e$} de $\mathfrak{g}$ tels que $(x,e)$ satisfassent aux
conditions de (i) forment un ouvert dense
\add{$\mathfrak{g}(2)^{*}$} de $\mathfrak{g}(2)$ (savoir l'ouvert des points
$e$ tels que $\operatorname{ad}(e)\colon\mathfrak{g}(0)\to\mathfrak{g}(2)$
soit surjectif), \add{Celui-ci} est un \struck{torseur} \add{espace
homogène} sous \add{$L=$} $\operatorname{Cent}(x)=\operatorname{Cent}(i)$,
qui est un groupe réductif.
\note{au-dessus de « éléments » un « $\neq 0$ » est biffé.}

(iv) Pour tout \add{élément de Dynkin $x$ de $\mathfrak{g}$,
$x=\operatorname{Lie}(i)(1)$,} \struck{couple $(x,e)$ satisfaisant les
conditions de (i)} considérons les sous-groupes $P$, $U$, \add{$L$} de $G$
définis par les conditions
\[
  \operatorname{Lie}P=\sum_{k\geqslant 0}\mathfrak{g}(k),\qquad
  \operatorname{Lie}U=\sum_{k\geqslant 1}\mathfrak{g}(k),\qquad
  \operatorname{Lie}L=\mathfrak{g}(0).
\]
Alors $P$ est un sous-groupe parabolique de $G$, $U$ est son radical
unipotent, et $L=\operatorname{Cent}(x)=\operatorname{Cent}(i)$ est un
sous-groupe de Lévi de $G$. Si $u$ \add{$=\exp(e)$} est un élément unipotent
de $G$ tel que $(i,u)$ soit un couple de Dynkin (i.e.\ induit par un
$f\colon\mathrm{Sl}(2,k)\to G$), $P$ (donc aussi $U$) est défini en termes de
$u$ (i.e.\ deux $x$ associés au même $u$ donnent le même $P$, cf.\ (ii))\,; il
en est de même du sous-groupe $H=i(\underline{G}_m)\cdot U$ de $G$ (NB
$i(\underline{G}_m)$ normalise $U$, et $H$ est un produit semi-direct), et de
$i\ \operatorname{mod}U$ i.e.\ de $\underline{G}_m\to H/U$ induit par $i$. De
plus, l'application
\[
  g\longmapsto\operatorname{ad}(g).x
\]
est un isomorphisme de $U$ avec $x+\operatorname{Fil}^{1}\mathfrak{g}$
\add{$=x+\operatorname{Lie}U$} (où on pose $\operatorname{Fil}^{k}
\mathfrak{g}=\sum_{j\geqslant k}\mathfrak{g}(j)$), et l'application
\[
  g\longmapsto\operatorname{ad}(g).e
\]
est une application surjective de $U$ sur $e+\operatorname{Fil}^{3}
(\mathfrak{g})$, ou de $P$ sur $\mathfrak{g}(2)^{*}+\operatorname{Fil}^{3}
(\mathfrak{g})$, où $\mathfrak{g}(2)^{*}$ est l'ouvert de $\mathfrak{g}(2)$
envisagé dans (iii).
\note{la frappe porte « $k\geqslant 0$ » corrigé à la main en
« $k\geqslant 1$ » sous le signe somme de $\operatorname{Lie}U$.}

\page{17}
\textbf{Interprétations en termes de filtrations du foncteur oubli} sur la
catégorie des représentations linéaires de dimension finie de $G$. (Il s'agit
de filtrations compatibles au produit tensoriel \add{et « exactes », cette
dernière cond.\ aut.\ si car.\ $k=0$}). Supposons fixé une classe \add{$c$}
de \add{conjugaison, rationelle sur $k$, de} telles filtrations (NB elle
n'admet pas nécessairement un représentant rat$/k$)\,; il revient au même de
se donner, pour le Chevalley $G_0$ associé à $G$, un élément de
$\operatorname{Hom}(\underline{G}_m,T_0)/W=M/W$
($T_0=M\otimes\underline{G}_m$ est le tore maximal type de $G_0$) invariant
par le sous-groupe \add{$\Gamma$} de $\operatorname{Autext}G_0=E_0$ image de
$\operatorname{Gal}(\bar k/k)$, ou encore un élément de $M$ qui soit dans la
chambre de Weyl de $M\otimes_{\mathbf{Z}}\mathbf{R}$ et qui soit stable par
$\Gamma$. On a alors le schéma $\underline{\varnothing}_c$ des filtrations de
classe $c$, qui est un espace homogène sous $G$, de stabilisateur des points
des paraboliques\,; \add{c'est un sous-schéma ouvert de
$\underline{\mathrm{Par}}(G)$, orbite rationnelle sous $G$, soit $c'$.} On
peut regarder aussi le schéma $\mathrm{ST}_c$ des homomorphismes
$i\colon\underline{G}_m\to G$ de classe $c$, qui est donc une orbite
rationelle de $G$ opérant sur $\underline{\mathrm{Homgr}}(\underline{G}_m,G)$.
Il est canoniquement isomorphe au schéma des couples de Lévi $(P,L)$ tels que
$P$ soit de classe $c'$ \add{ou encore $\underline{\varnothing}$}. (En
d'autres termes, grâce au choix d'un $c$, pour $P$ de classe $c'$ donné, il
revient au même de se donner un $i\colon\underline{G}_m\to P$ de la classe
définie par $c$, ou un sous-groupe de Lévi de $P$, à $i$ étant associé
$\operatorname{Cent}(i)$.) On a deux morphismes de schémas canoniquement
isomorphes
\[
  \mathrm{ST}_c\to\underline{\varnothing}_c,\qquad
  (\mathrm{CL})_{c'}\to\mathrm{Par}_{c'},
\]
le premier en termes de théorie des filtrations du foncteur fibres, le
deuxième en termes de groupes paraboliques.
\note{les passages marqués comme ajouts sont à la main ; celui qui suit
« paraboliques », interlinéaire, se lit mal (« c'est », « soit $c'$ »).}

Supposons maintenant que $k$ soit de car.\ nulle (j'ignore dans quelle mesure
c'est vraiment nécessaire). On peut alors considérer le Schéma $\mathrm{Dyn}_c$
des homomorphismes $f\colon\mathrm{Sl}(2,k)\to G$ tels que $f i_0$ soit de
classe $c$. Supposons-le non vide \add{et non réduit à un point,} i.e.\ $c$ une
classe de Dynkin \add{$\neq 1$}. Alors $\mathrm{Dyn}_c$ est un
\struck{torseur} \add{espace homogène} sous $G$ \add{de stabilisateurs $Z_0$
(il faudrait regarder la structure de ces $Z_0$\,; dimension, semi-simplicité).}
\struck{(en fait, il est can.\ is.\ au torseur
$\underline{\mathrm{Isom}}(G_1,G)$, où $G_1$ est le groupe quasi-déployé
correspondant à $G$).} Soit $\mathrm{Un}_c$ le schéma des éléments unipotents
\add{de $G$} dont la classe corresponde à $c$ par la correspondance

\page{18}
de Kostant. On a donc des morphismes $f\mapsto f\circ i_0$ et $f\mapsto
f(u_0)$ de $\mathrm{Dyn}_c$ dans $\mathrm{ST}_c$ resp.\ $\mathrm{Un}_c$.
D'autre part, Deligne nous fournit $\mathrm{Un}_c\to
\underline{\varnothing}_c$, et celui-ci peut aussi être défini par la
commutativité du diagramme \add{d'esp.\ homogènes}

\begin{tikzcd}
\mathrm{ST}_c\simeq(\mathrm{CL})_{c'} \arrow[d] & \mathrm{Dyn}_c \arrow[l] \arrow[d] \\
\underline{\varnothing}_c\simeq\mathrm{Par}_{c'} & \mathrm{Un}_c \arrow[l]
\end{tikzcd}

\begin{tikzcd}
L \arrow[d] & Z_0 \arrow[l] \arrow[d] \\
P & \mathfrak{z}(e) \arrow[l]
\end{tikzcd}

\note{ce second carré est dessiné à la main, à droite du premier, sous
« d'esp.\ homogènes » : il en donne, semble-t-il, les stabilisateurs. La
lettre notée $\mathfrak{z}$ est un 3 bouclé, et $Z_0$ récrit sur une autre
lettre.}

Pour qu'il existe un point rationnel de
$\underline{\varnothing}_c\simeq\mathrm{Par}_{c'}$, il faut (et suffit) qu'il
en existe un de $\mathrm{Dyn}_c$ \add{(i.e., si $Z_0=e$, que $G$ soit
quasi-déployé)}\,: \struck{i.e.\ que $G$ soit quasi-déployé\,;} en effet, $P$
étant choisi, on peut en choisir un sous-groupe de Lévi $L$, \add{ou ce qui
revient au même} un $i\colon\underline{G}_m\to G$ dans la classe $c$,
définissant $P$, et alors les $u=\exp(e)$ qui s'adaptent à $i$ forment un
ouvert dense de $\mathfrak{g}(2)$, qui a donc un point rationnel. On voit donc
que pour certaines classes de paraboliques (qu'on pourrait appeler les
paraboliques de \add{Dynkin-}Kostant) l'existence d'un parabolique de cette
classe rationnel sur $k$ implique déjà que $G$ est quasi-déployé. (NB Quand
$G$ n'est pas une forme intérieure, il faut faire attention cependant de
vérifier d'abord que la classe $c'$ de $P$ provient bien d'un $c$ de Dynkin
\emph{invariant par} $\Gamma$, et il n'est pas clair que cette propriété
\add{de $c$} résulte du fait que $c'$ est invariante par $\Gamma$, faute de
savoir si $c\mapsto c'$ est injectif sur les classes de Dynkin). (On notera
aussi que l'argument donné marche également sur un anneau semi-local (de car.\
nulle jusqu'à nouvel ordre).)
\marginal{Il n'est en effet \uncertain{pas} \ill{} \uncertain{injectif} pour
$A_2$, $A_4$ (\uncertain{cela vaut} pour $A_3$)}
Ceci montre l'intérêt, pour les questions arithmétiques, de déterminer
exactement les \add{classes de} sous-groupes paraboliques de Dynkin-Kostant
des divers groupes simples, et de voir si leur stabilisateur dans
$\operatorname{Autext}(G_0)$ est le même que celui des $c$ de Dynkin
correspondants, et pas str.\ plus grand.

On notera cependant que si une filtration du foncteur oubli peut être définie
par un élément unipotent, la classe de conjugaison de ce der-

\page{19}
nier (ou, ce qui revient au même, $c$) est bien déterminée. Les unipotents qui
définissent ladite filtration, en termes d'un $i$ qui la définit, sont
\add{ceux} de la forme $\exp(e)$, où $e$ parcourt
$\mathfrak{g}(2)^{*}+\operatorname{Fil}^{3}(\mathfrak{g})$. Ils forment une
classe de conjugaison d'éléments de $U$ \emph{par} $P$ (NB mais pas par $U$),
les classes de conjugaison \emph{par} $U$ s'identifiant aux ensembles de la
forme $(\exp(e)\,|\,e\in a+\operatorname{Fil}^{3}(\mathfrak{g}))$, où
$a\in\mathfrak{g}(2)^{*}$ est un élément quelconque.

\textbf{Cas principal}. Pour qu'on ait \ill{}$(1)=0$, il faut et suffit que
la classe d'homomorphismes $f\colon\mathrm{Sl}(2,k)\to G$ envisagée soit
« principale » au sens de Dynkin-Kostant. Alors le parabolique associé $P$ est
un Borel (la réciproque n'étant pas vraie \struck{sans doute}, cf table de
Dynkin), et le groupe $Z_0$ est le groupe unité (i.e.\ $\mathrm{Dyn}_c$ est un
torseur sous $G$), la réciproque n'étant pas vraie sans doute.
\marginal{\emph{ex} Cas $A_2$, $A_4$\,: les paraboliques associés aux
unipotents « sous-réguliers » sont des Borels\,!}
\note{devant « $(1)=0$ » la frappe laisse un blanc, que la main n'a pas
rempli : on attend $\mathfrak{g}(1)$. Le reste du feuillet est blanc.}
\end{document}
